\documentclass[10pt,a4paper]{article}
\usepackage[spanish,es-nodecimaldot]{babel}
\usepackage[utf8]{inputenc}
\usepackage[T1]{fontenc}
\usepackage{lmodern}
\usepackage[top=1.1cm,bottom=1.1cm,left=1.5cm,right=1.5cm]{geometry}
\usepackage{amsmath,amssymb}
\usepackage[table]{xcolor}
\usepackage{enumitem}
\usepackage[normalem]{ulem}

\setlength{\parindent}{0pt}
\setlength{\parskip}{2pt}
\pagestyle{empty}

\AtBeginDocument{%
  \setlength{\abovedisplayskip}{3pt}%
  \setlength{\belowdisplayskip}{3pt}%
  \setlength{\abovedisplayshortskip}{2pt}%
  \setlength{\belowdisplayshortskip}{2pt}%
}

\begin{document}

{\large\bfseries\uline{Pruebas de aleatoreidad}} --- Prueba de corridas
\vspace{4pt}

Sirven para comprobar si el orden de una muestra o proceso puede considerarse \textbf{aleatorio}. Se analiza si los valores se suceden sin un patrón sistemático.\\
Una \textbf{corrida} o \textbf{racha} es una secuencia de una o más letras iguales consecutivas.
\[
a\,a \mid b\,b \mid a \mid b\,b\,b \qquad \longrightarrow \quad \textbf{4 corridas}
\]

\textbf{\large Hipotesis}

H0:el arreglo o secuencia es aleatorio

H1:el arreglo no es aleatorio $\Rightarrow$prueba bilateral\\
H1:hay demasiadas corridas $\Rightarrow$prueba unilateral derecha\\
H1:hay muy pocas corridas $\Rightarrow$prueba unilateral izquierda

\textbf{\large Pasos}

\begin{minipage}[t]{0.55\textwidth}
\textbf{1. Construir la cadena de letras}

-Si los datos son \textbf{binarios} ---éxito/fracaso, 1/0, caro/barato---, se construye directamente una cadena de letras a y b.\\
-Si los datos son \textbf{numéricos}:
\begin{enumerate}[leftmargin=18pt,itemsep=0pt,topsep=1pt,parsep=0pt]
  \item Calcular la mediana.
  \item Asignar a a cada observación superior a la mediana.
  \item Asignar b a cada observación inferior a la mediana.
  \item Los valores iguales a la mediana se omiten.
  \item Mantener el orden original de los datos.
\end{enumerate}

\textbf{2. Contar las cantidades}

n1=cantidad de letras a\\
n2=cantidad de letras b\\
u=cantidad total de corridas

\textbf{3. Calcular los parámetros de la distribución de corridas}
\[
\mu_u = \frac{2n_1 n_2}{n_1+n_2} + 1
\]
\[
\sigma_u = \sqrt{\frac{2n_1 n_2\,(2n_1 n_2 - n_1 - n_2)}{(n_1+n_2)^2\,(n_1+n_2-1)}}
\]
\end{minipage}\hfill
\begin{minipage}[t]{0.42\textwidth}
\textbf{4. Aproximacion normal}

si n1$>$10 y n2$>$10 la dist. del num de corridas \textbf{u} se aproxima a una dist. normal, aplico estadistico z
\[
Z_{\text{calc}} = \frac{u - \mu_u}{\sigma_u}
\]

\textbf{Interpretaciones}
\begin{itemize}[leftmargin=12pt,itemsep=1pt,topsep=1pt,parsep=0pt]
  \item {\small $u$ grande: demasiados cambios entre $a$ y $b$; alternancia excesiva.}
  \item {\small $u$ chico: pocas alternancias; los valores iguales tienden a agruparse.}
  \item {\small $u$ cercano a $\mu_u$: comportamiento compatible con aleatoriedad.}
\end{itemize}

\vspace{4pt}
\textbf{Se usa en control de calidad para dar aleatoreidad a procesos}
\end{minipage}

\end{document}
